% Part: free-logic
% Chapter: natural-deduction
% Section: common-rules

\documentclass[../../../include/open-logic-section]{subfiles}

\begin{document}

\olfileid{frl}{dis}{com}

\olsection{Comparisons Between First-Order Logic and Free Logics}

\begin{center}
Comparing validity across first-order logic with identity (\Log{FOL}) 
and free logics: negative free logic (\Log{NFL}) and positive free 
logic (\Log{PFL}).
\smallskip
\begin{tabular}[t]{lp{.33\textwidth}ccc}
    \hline
    !!^{sentence} & English trans. & \multicolumn{3}{c}{Valid in:} \\
        &       & \Log{FOL} & \Log{NFL} & \Log{PFL} \\
    \hline 
    \multicolumn{2}{l}{Classical vs free quantification} \\
    $\lforall[x][!A]\lif\Subst{!A}{\Obj{a}}{x}$ 
        & If everything is $!A$, $\Obj{a}$ is $!A$ 
        & $\surd$ & $\times$ & $\times$ \\
    $\lforall[x][!A]\lif(\lfrexists\Obj{a}\lif\Subst{!A}{\Obj{a}}{x})$ 
        & If everything is $!A$, then if $\Obj{a}$ exists $\Obj{a}$ is $!A$
        & $\surd$ & $\surd$ & $\surd$ \\
    \multicolumn{2}{l}{Classical vs free existence assumptions} \\    
    $\lforall[x][\lfrexists x]$ 
        & Everything exists 
        & $\surd$ & $\surd$ & $\surd$ \\
    $\lfrexists \Obj{a}$ 
        & $\Obj{a}$ exists 
        & $\surd$ & $\times$ & $\times$ \\
    \multicolumn{2}{l}{Identity: \Log{PFL} like \Log{FOL}, \Log{NFL} not} \\
    $\lforall[x][\eq[x][x]]$
        & Everything is itself
        & $\surd$ & $\surd$ & $\surd$ \\
    $\eq[\Obj{a}][\Obj{a}]$
        & $\Obj{a}$ is $\Obj{a}$
        & $\surd$ & $\times$ & $\surd$ \\
    \multicolumn{2}{l}{Atomic predication: \Log{NFL} like \Log{FOL}, \Log{PFL} not} \\
    $\Atom{\Obj{P}}{\Obj{a}}\lif\lfrexists\Obj{a}$ 
        & If $\Obj{a}$ is $\Obj{P}$, then $\Obj{a}$ exists 
        & $\surd$ & $\surd$ & $\times$ \\
    \multicolumn{2}{l}{Negated atomic predication: classical vs free} \\
    $\lnot \Atom{\Obj{P}}{\Obj{a}}\lif\lfrexists\Obj{a}$ 
        & If it's not the case that $\Obj{a}$ is $\Obj{P}$, then $\Obj{a}$ exists 
        & $\surd$ & $\times$ & $\times$ \\
    \hline
\end{tabular}
\smallskip
\end{center}

\emph{Note on the existence predicate $\lfrexists$}. In free logics
it is equivalent to $\lexists[v][\eq[\ldots][v]]$. In first-order
logic with identity we haven't introduced it, but it would be
equivalent to that to. Thus were first-order logic is concerned,
simply treat $\lfrexists c$ as an abbreviation of
 $\lexists[v][\leq[c][v]]$.
Thus:
\begin{itemize}
    \item $\lfrexists \Obj{a}$ is 
    equivalent to $\lexists[x][\eq[\Obj{a}][x]]$
    \item $\lforall[x][\lfrexists x]$ is 
    equivalent $\lforall[x][\lexists[y][\eq[x][y]]]$
    \item $\lfrexists\Obj{a}\lif\Subst{!A}{\Obj{a}}{x}$ is 
    equivalent $\lexists[y][\eq[\Obj{a}][y]]\lif\Subst{!A}{\Obj{a}}{x}$
\end{itemize}

\end{document}